\section{Whole traces and reflected high indices}\label{ms:fourier}
The following comparisons concern the actual top-coefficient trace.
They retain all its Fourier aliases, including the alternating signs
at even levels. Put $\epsilon=(-1)^{N+1}$.

\begin{lemma}[The finite Fourier trace]\label{ms:Fourier}
If $F(-u-u^{-1})=\sum a_k u^k$, then
\begin{equation}\label{ms:Fourier-trace}
 \phi_s(F)=N^2\sum_{l\ge1}\epsilon^{l+1}
    (a_{lN-1}-a_{lN+1})U_{l-1}(\cos2\pi s).
\end{equation}
Only finitely many summands occur. The formula holds at every
collision with the full normal functional.
\end{lemma}
\begin{proof}
At generic $s$ choose $u_j^N=\epsilon e^{2\pi\ii s}$ and put
$z_j=-u_j-u_j^{-1}$. These are the distinct roots of $P_s$, and
\[
 P_s'(z_j)=\frac{2\ii N\sin2\pi s}{u_j-u_j^{-1}},\qquad
 \phi_s(F)=\frac{\epsilon N}{2\ii\sin2\pi s}
       \sum_j(u_j-u_j^{-1})F(-u_j-u_j^{-1}).
\]
The root sum selects multiples of $N$. Pairing opposite exponents
gives \eqref{ms:Fourier-trace}, since $a_{-k}=a_k$. Both sides are
polynomials in $2\cos2\pi s$: the left by monic reduction, the right
by the finite Chebyshev sum. Polynomial continuation proves the
collision assertion, including all first normals.
\end{proof}

For low $n$, the exact coefficients are positive:
\begin{equation}\label{ms:low-coefficients}
 [u^k]e_n(-u-u^{-1})=
 \frac{S_{2n}}{S_n^2S_{n+k}S_{n-k}},\qquad |k|\le n.
\end{equation}
Indeed the numerator is
$u^{-n}\prod_{a<n}(1+q^{2a+1}u)(1+q^{-2a-1}u)$.
Expand the consecutive geometric factors by the finite binomial
identity, proved by multiplying by its last factor. The phases
cancel and give \eqref{ms:low-coefficients}. Every factorial index
is less than $N$. Symmetry is explicit; the successive ratio on the
nonnegative side is
$\sin((n-k)\pi/N)/\sin((n+k+1)\pi/N)\le1$.
Thus these sequences are symmetric and unimodal. Their convolutions
are too: write a symmetric unimodal sequence as a nonnegative sum
of centered interval indicators.

We will also use an elementary quantitative tail fact. If $a$ is
the convolution of $B$ nonnegative symmetric unimodal sequences
supported in $[-m,m]$, then
\begin{equation}\label{ms:interval-tail}
 a_k\le(2m+1)^B(a_k-a_{k+2}),\qquad k\ge m.
\end{equation}
To prove it after the interval decomposition, let $C$ be the
convolution of integer interval indicators, with largest radius $r$.
Write $C=\mathbf1_{[-r,r]}*H$. For $k\ge r$,
$C(k)-C(k+1)=H(k-r)-H(k+r+1)$. If $C(k)>0$, the first
argument is in the support. If the second is outside, the difference
is positive. Otherwise induction makes its last step strict, since
$k+r$ is at least every radius in $H$; unimodality handles the other
steps. The one-factor case is immediate. The positive integer
difference is at least one, while $C(k)\le(2m+1)^B$. Sum the
interval decompositions to obtain \eqref{ms:interval-tail}.

\begin{proposition}[Same-level low trace comparison]\label{ms:low-trace}
For fixed $B$ and any bounded set of complex $s$,
\begin{equation}\label{ms:low-trace-bound}
 |\phi_s(e_{n_1}\cdots e_{n_B})|
       \le C N^A\phi_0(e_{n_1}\cdots e_{n_B}),\qquad n_i\le m.
\end{equation}
A zero on the right forces the left to vanish. Both parities are
included, and all factors on the right are at the original level.
\end{proposition}
\begin{proof}
Write $D=\sum n_i$ and $a_k$ for the product coefficients.
All $d_l=a_{lN-1}-a_{lN+1}$ are nonnegative. The number of aliases
in \eqref{ms:Fourier-trace} is at most $B$, and
\[
 |U_{l-1}(\cos2\pi s)|\le l e^{2\pi(l-1)|\Im s|}.
\]
For odd $N$ this proves the result directly, since all complete
alias signs are positive. If $D<N-1$, both traces vanish by degree.
It remains to consider even $N\ge4$, $D\ge N-1$.

Put $L=N-1$, an odd level; its low bound is the same $m=(N-2)/2$.
Concavity of sine and reflection about $\pi/2$ give, for
$0\le r\le N-2$,
\begin{equation}\label{ms:adjacent-products}
 \left(\frac N{N-1}\right)^r\frac{N-1-r}{N-1}
 \le\frac{S_r^{(L)}}{S_r^{(N)}}
 \le\left(\frac N{N-1}\right)^r.
\end{equation}
For the lower estimate multiply
$\sin(\pi a/(N-1))/\sin(\pi a/N)
\ge N(N-1-a)/((N-1)(N-a))$; the product telescopes.
The upper estimate follows from decreasing $\sin t/t$.
All these ratios and their inverses are polynomial in $N$.
Changing an allowed factorial index by one also costs at most a
polynomial, since $4/N\le2\sin(\pi a/N)\le2$.
Consequently \eqref{ms:low-coefficients} and convolution give
$a_k^{(N)}\le N^A a_k^{(L)}$ at every index, with exactly the
same support. Unimodality and \eqref{ms:interval-tail}, now at
$k=lL-1\ge m$, imply
\[
 d_l^{(N)}\le N^A a_{lN-1}^{(L)}
 \le N^A a_{lL-1}^{(L)}
 \le N^{A+B}(a_{lL-1}^{(L)}-a_{lL+1}^{(L)}).
\]
The odd complete trace therefore bounds the absolute even trace
up to a polynomial.

We must return to the original level. The complete multiplication
law, \cite[Theorem 3.3]{Newton}, is positive on
\begin{equation}\label{ms:level-gate}
 \mu_K(i,j;r)=
 \frac{S_{i+j}^{(K)}S_r^{(K)}}
 {S_i^{(K)}S_j^{(K)}S_{r-i}^{(K)}S_{r-j}^{(K)}S_{i+j-r}^{(K)}},
 \quad i+j<K,\quad\max(i,j)\le r\le i+j;
\end{equation}
it is zero otherwise. A successive odd-level trace path starts at
$r_1=n_1$, ends at $r_B=N-2$, and has nonnegative gate slacks
$r_{j-1}+n_j-r_j$. Their sum is $D-(N-2)\ge1$. At its first
positive slack raise that output and every subsequent running
output by one. At the first gate the slack permits the increase;
at later gates the input sum and output both increase by one, so
all support inequalities persist and input sums remain at most
$N-1$. The new path closes at $N-1$, the even trace target.
Every image has at most $B-1$ preimages, since specifying the first
changed gate recovers the path. Each factorial index changes by at
most one. Equation \eqref{ms:adjacent-products} thus bounds each
old positive path weight by a polynomial times its image weight.
Summing proves $\phi_{L,0}\le N^A\phi_{N,0}$, completing the proof.
This level comparison is only an estimate, not an identification of
invariants at different levels.
\end{proof}

\begin{proposition}[A reflected high trace]\label{ms:reflected}
Let $h>m$, $r=d-h$ and $F=\prod_{j=1}^B e_{n_j}$, $n_j\le m$.
For fixed $B$ and bounded $s$,
\begin{equation}\label{ms:reflection-bound}
 |(\phi_s-\phi_0)(e_hF)|
 \le C N^A|\Delta_s|\,\phi_0(e_rF),
 \qquad \Delta_s=2\cos2\pi s-2.
\end{equation}
If $n_j/N\in[\delta,1/2-\delta]$ and $r/N\ge\delta>0$,
the right side can be multiplied by $e^{-\kappa N}$, with
$\kappa>0$ independent of the meridian compact.
\end{proposition}
\begin{proof}
Set $a=N-h$, $v=2h-N\ge0$. Removing one complete cycle of
geometric factors from the Newton numerator gives
\[
 \prod_{j<h}(1+q^{2j+1}u)(1+q^{-2j-1}u)
 =(1-\epsilon u^N)\prod_{j=0}^{v-1}(1-q^{2j-v+1}u).
\]
The last product has coefficient $(-1)^l S_v/(S_lS_{v-l})$.
Thus the Laurent coefficients $A_k$ of $e_h(-u-u^{-1})$ are
zero for $|k|<a$, and otherwise
\begin{equation}\label{ms:high-coefficients}
 A_k=(-1)^{h-|k|}
       \frac{S_{2h-N}}{S_h^2S_{|k|-a}S_{h-|k|}},
       \qquad a\le|k|\le h.
\end{equation}
At even $h=N/2$ the only coefficients are
$A_h=A_{-h}=S_h^{-2}$. By \eqref{ms:basis-norm} and polynomial
conditioning, the interval norm of this high polynomial is
polynomial; its Laurent coefficients are cosine coefficients, so
$\sum|A_k|\le N^A$.

Write $F(-u-u^{-1})=\sum f_k u^k$. Its coefficients are positive
and unimodal. Alias one in \eqref{ms:Fourier-trace} is independent
of $s$ and cancels in the difference. All others have frequencies
at least $2N-1$, so convolution gives a bound by
$N^A f_{2N-1-h}=N^A f_{N+r}$. Moreover
$U_{l-1}(\cos2\pi s)-l$ is divisible by $\Delta_s$ with bounded
quotient for the fixed finite set of aliases. This proves
\begin{equation}\label{ms:high-tail}
 |(\phi_s-\phi_0)(e_hF)|\le C N^A|\Delta_s|f_{N+r}.
\end{equation}
If $C=e_rF=\sum c_k u^k$, the endpoint coefficient of $e_r$ is
$S_r^{-2}$, and $S_r^2\le N$ on the low range. Hence
\[
 c_{N-1}\ge S_r^{-2}f_{N-1-r}\ge N^{-1}f_{N+r}.
\]
The tail lemma bounds $c_{N-1}$ by a polynomial times
$d_1=c_{N-1}-c_{N+1}$. To bound this first difference by the
complete trace even when its aliases alternate, form
$Q(X)=\sum_l\epsilon^{l+1}d_lU_{l-1}(X)$.
Its degree is bounded by $B$. Evaluate at $B+1$ fixed distinct
values $X_j=\cos2\pi s_j$, $0\le s_j<1/8$. The inverse
evaluation matrix is fixed, so
$d_1\le C\max_j|Q(X_j)|\le N^A\phi_0(C)$ by
Proposition~\ref{ms:low-trace}. This proves
\eqref{ms:reflection-bound} at the same original level.

For strictness retain the two different frequencies above. Define
\[
 b_u(t)=J(2u)-2J(u)-J(u+t)-J(u-t),\quad |t|\le u.
\]
On $u\in[\delta,1/2-\delta]$,
\[
 b_u''(t)=-\frac{\pi\sin(2\pi u)}
                 {\sin\pi(u+t)\sin\pi(u-t)}\le-\lambda_\delta<0.
\]
The constrained convolution action
$\mathcal B(t)=\max_{\sum t_j=t}\sum b_{u_j}(t_j)$
is even and strongly concave with constant $\lambda_\delta/B$,
by Cauchy--Schwarz on the coordinate differences of two maximizers.
It follows that
\begin{equation}\label{ms:convolution-drop}
 \mathcal B(t_2)-\mathcal B(t_1)
 \le-\frac{\lambda_\delta}{2B}(t_2^2-t_1^2),\quad 0\le t_1\le t_2.
\end{equation}
Uniform sine-factorial estimates give
$\log f_k=N\mathcal B(k/N)+O(\log N)$ on the support.
For the lower bound round a maximizing tuple, preserving its integer
sum: floor all coordinates and raise the required number with
nonzero fractional parts. The integer endpoints remain admissible.
The modulus $|J(x+h)-J(x)|\le C|h|\log(e/|h|)$ controls even
endpoint rounding; the upper bound counts polynomially many tuples.
Apply \eqref{ms:convolution-drop} at
$(N-1-r)/N$ and $(N+r)/N$. When the latter coefficient is nonzero,
\[
 f_{N+r}\le N^A e^{-2\lambda_\delta r/B+O(1)}f_{N-1-r}.
\]
Since $r\ge\delta N$, this supplies the asserted strict factor
in the preceding comparison. If it is zero, the difference is zero.
\end{proof}

\begin{corollary}[Regular low windows]\label{ms:regular-trace}
If all $n_j/N\in[\delta,1/2-\delta]$ and $\sum n_j/N>1+\delta$,
then $\phi_s(F)/\phi_0(F)=1+O(e^{-\kappa N})$ uniformly on any
fixed meridian compact, with fixed derivatives. For at most four
low factors the traces agree exactly, without a window assumption.
\end{corollary}
\begin{proof}
Use the same convolution action at $t_1=1-1/N$ and
$t_2=l-1/N$, $l\ge2$. Their squared difference is at least two.
The tail lemma gives $d_1\ge N^{-B}a_{N-1}>0$, so the sum of all
other $d_l$ is exponentially smaller. Formula
\eqref{ms:Fourier-trace} at $s$ and at zero proves the assertion,
including the even signs. With four low factors the raw degree is
at most $4m\le2N-2$, so there is no second alias. Cauchy's estimate
on a larger fixed meridian compact proves the derivative statement.
\end{proof}
